\section*{Abstract}
\addcontentsline{toc}{section}{Abstract}

If a computer could predict how a cell's genes will respond to a drug, compounds could be screened
without running the experiment. Recent work makes this look within reach: a cell can be written down
as a sentence --- the names of its active genes, listed from most active to least --- which lets a
language model of the kind used for text be trained to write the ``sentence'' of a treated cell given
an untreated one and the name of a drug. This thesis trains such a model on the largest available
atlas of drug-treated single cells and asks a simple question: does it learn anything about the
drugs?

It does not, and establishing that first required fixing the way such models are graded. The standard
score compares predicted against observed gene changes. A predictor that ignores its inputs entirely,
guessing that every gene sits in the middle of the pack, scores $0.9999$ on it --- higher than two
repeats of the same real experiment score against each other. Less information earns a better grade,
which is the signature of a broken measurement. Only one metric examined here survives a formal
calibration test.

Graded on that one, the model does no better than chance, and no better than when the drug in its
prompt is swapped for another at random. The information is not missing: the model's internal
activity identifies the drug ever more clearly through its layers, while the part of the network that
produces the output attends to it almost not at all. The cause lies in how the target is written.
Genes are listed by activity, the most active genes are largely the same whatever drug is applied,
and so $83\%$ of the words in the target sentence are identical for any two drugs --- leaving almost
no training signal about the drug. Listing genes instead by \emph{how much the drug changed them}
raises that to $59\%$ and produces the first measurable drug effect in this work, one that carries
over to drug--cell-line combinations never seen in training.

The model is nonetheless beaten by a lookup table that simply recalls what the same drug did in other
cell lines, and decomposing the response explains why. About half of a drug's effect is a fixed
signature it carries everywhere, which the lookup reproduces almost perfectly. The other half depends
on the particular cell line --- and that half has no consistent structure, in that no cell line bends
different drugs in a consistent direction. Being unpredictable, it limits every method equally, so
for a drug already seen the lookup is as good as anything can be. The opening is drugs that have
\emph{not} been seen, where a lookup is impossible by construction: there a drug's protein targets
and mechanism carry usable information, while its chemical structure does not. But the model ignores
the mechanism it is already told, so the obstacle there is again what the network attends to rather
than what it is given.

\vspace{0.6em}
\noindent\textbf{Keywords:} single-cell transcriptomics, drug perturbation prediction, language
models, benchmark calibration, mechanistic interpretability.
